% Apêndice: figuras complementares. Não contam no limite de 10 a 15 páginas do corpo.
\appendix
\section{Figuras complementares}\label{sec:apendice_figuras}

As figuras deste apêndice detalham partes do sistema resumidas na \autoref{sec:siras}.

\begin{figure}[htb]
\centering
\figura{F2-arquitetura}{\columnwidth}
\caption{Camadas do SIRAS e dependências entre os pacotes. Nenhuma camada abaixo da interface
depende do Flask.}
\label{fig:arquitetura}
\end{figure}

\begin{figure*}[htb]
\centering
\figura{F6-interpretacao}{\textwidth}
\caption{Interpretação de fósforo pela classe de teor de argila e de potássio pela CTC a pH
7,0, para o grupo de exigência das culturas de grãos (Manual 2016, Tab. 6.4 e 6.9). O ponto
marca a Gleba~4 do Anexo~1.}
\label{fig:interpretacao}
\end{figure*}

\begin{figure*}[htb]
\centering
\figura{F7-adubacao}{\textwidth}
\caption{Os dois caminhos da adubação: o algoritmo dos grãos e a dose por classe dos demais
grupos, parametrizada pelas variáveis de cada um.}
\label{fig:adubacao}
\end{figure*}

\begin{figure*}[htb]
\centering
\figura{F8-formulacao}{\textwidth}
\caption{Os três caminhos da formulação do adubo e a escolha da dose pelos pontos de quebra,
ilustrados com a saída do motor para a Gleba~4.}
\label{fig:formulacao}
\end{figure*}

\begin{figure}[htb]
\centering
\figura{F10-trilha}{\columnwidth}
\caption{Anatomia de um passo da trilha de inferência e sua apresentação no laudo.}
\label{fig:trilha}
\end{figure}

\begin{figure*}[htb]
\centering
\figura{F12-fluxo-de-uso}{\textwidth}
\caption{Fluxo de uso do SIRAS na perspectiva do usuário.}
\label{fig:fluxo}
\end{figure*}

\begin{figure}[htb]
\centering
\figura{F14-usabilidade}{\columnwidth}
\caption{Desenho da sessão de avaliação de usabilidade.}
\label{fig:sessao}
\end{figure}
